\documentclass[11pt,a4paper]{article}

\usepackage[T1]{fontenc}
\usepackage[utf8]{inputenc}
\usepackage[margin=2.5cm]{geometry}
\usepackage{amsmath,amssymb}
\usepackage{booktabs}
\usepackage{hyperref}
\usepackage{parskip}
\usepackage{listings}
\usepackage{xcolor}
\usepackage{float}

\lstset{
  basicstyle=\ttfamily\small,
  keywordstyle=\color{blue},
  commentstyle=\color{gray},
  backgroundcolor=\color{gray!8},
  frame=single,
  framerule=0.4pt,
  language=Python,
  breaklines=true,
  showstringspaces=false,
}

\hypersetup{colorlinks=true, linkcolor=blue, urlcolor=blue, citecolor=blue}

\title{\textbf{Project 7: Stream Analysis}\\[4pt]
       \large Algorithms for Massive Datasets\\
       Master in Computer Science - Università degli Studi di Milano, 2025/26}
\author{Gian Marco Bratzu}
\date{June 2026}

\begin{document}
\maketitle

\begin{abstract}
This report describes the implementation of three streaming algorithms applied to the
\textit{New York Times Articles \& Comments 2020} dataset.
The comment log is interpreted as a time-ordered stream, and three algorithms are
implemented from scratch: \textbf{Flajolet-Martin} (count of distinct users),
\textbf{Alon-Matias-Szegedy} (second moment of the section distribution), and a
\textbf{Bloom filter} (set-membership test for articles belonging to a given section).
Each estimator is compared against the computed exact value to measure the
approximation error. All three work in a single pass and use memory independent of
the stream length.
The full code is available at:
\url{https://github.com/Bratzu93/Algorithms-for-massive-datasets}.
\end{abstract}

\section{Introduction}

Streaming algorithms address the scenario in which data arrives as a continuous sequence
of elements that can be processed only once, without the ability to store the entire stream.
The fundamental constraints are one-pass processing and sublinear (ideally constant)
memory usage. In exchange, only approximate answers are provided, with provable
probabilistic guarantees.

This project applies three classical streaming techniques to user comments from
The New York Times, treating the sequence of comments, ordered by publication
timestamp as the stream.
The three problems addressed are:
\begin{enumerate}
  \item Counting \textbf{distinct users} (Flajolet-Martin), specifically the
        0\textsuperscript{th} frequency moment $F_0$.
  \item Estimating the \textbf{second moment} of the comment distribution across
        sections (Alon-Matias-Szegedy): how much does traffic concentrate on a few sections?
  \item \textbf{Bloom filter} for set membership: given an article ID, is it in Opinion?
\end{enumerate}

\section{Dataset}

\subsection*{Source}
New York Times Articles \& Comments (2020), Kaggle, identifier
\texttt{benjaminawd/new-york-times-articles-comments-2020} (CC BY-NC-SA 4.0, accessed June 2026)~\cite{kaggle}.

\subsection*{Files}
\begin{itemize}
  \item \texttt{nyt-articles-2020.csv}: 16,787 articles with fields including
        \texttt{uniqueID}, \texttt{section}, \texttt{headline}, \texttt{pub\_date}.
        The 42 distinct sections include U.S.\ (2,364 articles), Opinion (2,272),
        World (1,183), Arts (1,094), and New York (1,055).
  \item \texttt{nyt-comments-part0.csv}: first partition of the comments,
        approximately 300~MB and 500,000 rows. Used for the subsampled run.
  \item \texttt{nyt-comments-2020.csv}: single pre-merged file of approximately
        3~GB containing \textbf{4,986,461 comments} in total. Used for the full-scale run.
        Relevant fields (both files): \texttt{userID} (integer), \texttt{createDate} (timestamp),
        \texttt{articleID} (URI string).
\end{itemize}

\subsection*{Download}
The dataset is downloaded via the Kaggle CLI at notebook startup.
On Google Colab the credentials (\texttt{KAGGLE\_USERNAME}, \texttt{KAGGLE\_KEY})
are read from Colab Secrets; locally from \texttt{\textasciitilde/.kaggle/kaggle.json}.
The dataset is not committed to the repository.

\section{Data Organisation and Pre-processing}

\subsection*{Stream construction}
The stream is the sequence of comments sorted in ascending order by \texttt{createDate}, since the timestamp format is \texttt{YYYY-MM-DD HH:MM:SS}.

\subsection*{Section enrichment}
The comments file does not carry section information; it is added via a pandas
left join on the articles file, matching \texttt{comments.articleID} to
\texttt{articles.uniqueID}.
The articles file is loaded with only the two columns needed (\texttt{uniqueID},
\texttt{section}) and merged into the comments DataFrame in a single pass;
\texttt{uniqueID} is then dropped, leaving a \texttt{section} column alongside
each comment.
Manual inspection confirmed that every comment matches an article in the articles
file, so no \texttt{NaN} values appear in the \texttt{section} column after the join.

\subsection*{Subsampling}
A global boolean variable \texttt{SUBSAMPLE} controls data size at runtime.
\begin{itemize}
  \item \texttt{SUBSAMPLE = True} (default): loads \texttt{nyt-comments-part0.csv},
        sorts by timestamp, and retains the first 200,000 comments. This setting
        is designed for Google Colab and completes
        in a few minutes.
  \item \texttt{SUBSAMPLE = False}: loads \texttt{nyt-comments-2020.csv} directly and
        sorts the full 4,986,461-comment stream. Intended for local execution.
\end{itemize}

\subsection*{Stream simulation}
The stream is simulated by sorting the DataFrame by \texttt{createDate} and
iterating over the elements with a \texttt{for} loop, passing one item at a time
to each algorithm's \texttt{update} method, mirroring how a real stream would
trigger updates on arrival.
The full dataset is held in memory as input, but the streaming constraint applies
to the \emph{algorithm state}: FM stores $H$ integers and AMS stores $k$ counters,
both independent of the stream length.
The Bloom filter stores $m$ bits, where $m$ is computed from the number of elements
to insert $n$ and the target FPR; its memory therefore grows with the positive set
size, not with the stream length.
One acknowledged limitation is that AMS requires knowing $n$ before the stream
begins in order to sample positions uniformly; in a truly online setting this
would require a preliminary counting pass.

\section{Algorithms and Implementations}

\subsection{Flajolet-Martin (Distinct Elements)}

The goal is to estimate $F_0 = |\{x : f_x > 0\}|$, the number of distinct
\texttt{userID} values, in a single pass over the stream.

\subsubsection*{Algorithm}
For a hash function $h$ that maps elements uniformly over $\{0,1\}^b$, the probability
that $h(x)$ ends in at least $k$ trailing zero bits is $2^{-k}$. Consequently, if $R$
is the maximum number of trailing zeros observed across all elements in the stream,
then $\hat{F}_0 = 2^R$ is an estimate of $F_0$.

To reduce variance, $H$ independent hash functions are used.
Their $R_{\max}$ values are divided into $G$ groups; within each group the median is
taken (to suppress outliers arising from $h(x)=0$, which would give an infinite tail
length), and the final estimate is the mean of the per-group estimates:
\[
  \hat{F}_0 = \frac{1}{G}\sum_{g=1}^{G} 2^{\,\mathrm{median}(R^{(g)})}
\]

\subsubsection*{Initial approach: MD5 with sequential salts}

A first implementation derived $H$ hash functions from MD5 by varying a numeric salt:
\[
  h_i(x) = \mathrm{MD5}\!\left(\texttt{str}(x) + \texttt{"\_"} + \texttt{str}(\mathrm{seed}+i)\right)
  \bmod 2^{128}, \qquad i = 0, 1, \ldots, H-1
\]
The trailing zero count was read by converting the integer output to its binary
string representation and iterating over characters from the right.
This approach was replaced for two reasons.
First, \textbf{performance}: the MD5 version required approximately 30 seconds on
the subsample and 20 minutes on the full dataset, versus 3 seconds and 4 minutes
respectively for the universal hash implementation — a roughly 5--10$\times$ slowdown
attributable to string conversion and the character-by-character trailing-zero loop.
Second, \textbf{implementation}: the bit-trick \texttt{(h \& -h).bit\_length() - 1}
computes the trailing-zero count in $O(1)$ without any string parsing, and
hash functions of the form $(ax+b) \bmod p$ provide stronger theoretical
independence guarantees than MD5 with sequential salts.

\subsubsection*{Final implementation: universal hashing}
Each hash function is a member of the 2-universal family:
\[
  h_i(x) = (a_i \cdot x + b_i) \bmod p, \qquad
  p = 2^{61}-1 \text{ (Mersenne prime)}
\]
with $a_i \in [1, p-1]$ and $b_i \in [0, p-1]$ chosen independently at random.
Coefficients sampled from a space of $\sim 2^{61}$ values ensure statistical
independence across the $H$ hash functions, enabling genuine variance reduction.

\subsubsection*{Trailing zeros}
The number of trailing zero bits of $h$ is computed with the identity:
\begin{lstlisting}
(h & -h).bit_length() - 1
\end{lstlisting}
The expression \texttt{h \& -h} isolates the least significant set bit
(e.g.\ $12_{10} = 1100_2 \;\to\; 0100_2 = 4$); \texttt{bit\_length() - 1} gives its
position, which equals the trailing-zero count. I used \texttt{FM\_NUM\_HASHES = 64}
and \texttt{FM\_NUM\_GROUPS = 8}.

\subsubsection*{Results}

\begin{table}[H]
\centering
\begin{tabular}{r rr rr}
\toprule
& \multicolumn{2}{c}{MD5} & \multicolumn{2}{c}{Universal hash} \\
\cmidrule(lr){2-3} \cmidrule(lr){4-5}
$H$ & Estimate & Error & Estimate & Error \\
\midrule
  8 &  55\,296 &   8.7\% & 266\,240 & 423.3\% \\
 16 &  53\,745 &   5.6\% &  69\,426 &  36.5\% \\
 32 &  45\,347 &  10.9\% &  96\,281 &  89.2\% \\
 64 &  61\,028 &  20.0\% &  53\,042 &   4.3\% \\
128 &  54\,945 &   8.0\% &  48\,449 &   4.8\% \\
256 &  56\,144 &  10.4\% &  59\,041 &  16.0\% \\
512 &  57\,344 &  12.7\% &  63\,137 &  24.1\% \\
\bottomrule
\end{tabular}
\caption{FM on the subsample (200,000 comments, $G=8$ groups, \texttt{SEED=1},
  exact $F_0^* = 50{,}876$). Run time for $H=256$--$512$: MD5 2m\,40s, universal hash 33s.}
\label{tab:fm_subsample}
\end{table}

\begin{table}[H]
\centering
\begin{tabular}{r rr rr}
\toprule
& \multicolumn{2}{c}{MD5} & \multicolumn{2}{c}{Universal hash} \\
\cmidrule(lr){2-3} \cmidrule(lr){4-5}
$H$ & Estimate & Error & Estimate & Error \\
\midrule
  8 & 1\,933\,312 & 379.7\% & 663\,552 &  64.6\% \\
 16 &   992\,155 & 146.2\% & 534\,568 &  32.6\% \\
 32 &   409\,118 &   1.5\% & 465\,056 &  15.4\% \\
 64 &   501\,800 &  24.5\% & 540\,190 &  34.0\% \\
128 &   499\,471 &  23.9\% & 374\,021 &   7.2\% \\
\bottomrule
\end{tabular}
\caption{FM on the full dataset (4,986,461 comments, $G=8$ groups, \texttt{SEED=1},
  exact $F_0^* = 403{,}025$). Run time: MD5 20m\,10s, universal hash 4m\,20s.}
\label{tab:fm}
\end{table}

Universal hashing is approximately $5\times$ faster than MD5 in both settings.

This illustrates a key limitation of single-run evaluation: a single extreme
$R_{\max}$ value can shift a group median regardless of how many hash functions
are used, making the error non-monotone in $H$.
The formula gives the expected error over many independent runs, not a per-run guarantee.

The number of groups $G$ controls the confidence level: a larger $G$ reduces the
probability that all group medians are simultaneously wrong.
For robust estimates, averaging multiple independent runs (each with a different seed)
is more reliable than relying on a single run with a large $H$.

\subsection{Alon-Matias-Szegedy (Second Moment)}

The second frequency moment $F_2 = \sum_{j} m_j^2$ (where $m_j$ is the number of
comments in section $j$) measures how concentrated the distribution is across sections:
the more uneven, the higher $F_2$.

\subsubsection*{Algorithm}
Assume the stream length $n$ is known in advance. Sample $k$ positions
$p_1, \ldots, p_k$ uniformly at random from $[0, n)$.
For each sampled position $p_i$, maintain a variable $X_i$:
\begin{enumerate}
  \item When the stream reaches position $p_i$: record $\mathrm{elem}_i \leftarrow \mathrm{stream}[p_i]$ and set $v_i \leftarrow 1$.
  \item For each subsequent position $j > p_i$: if $\mathrm{stream}[j] = \mathrm{elem}_i$, increment $v_i$.
\end{enumerate}
The estimator for $F_2$ is:
\[
  \hat{F}_2 = \frac{1}{k}\sum_{i=1}^{k} n\cdot(2v_i - 1)
\]
It can be shown that $\mathbb{E}[n(2v_i-1)] = F_2$, so the average over all $k$
variables gives an unbiased estimate. The intuition is that frequent elements are
more likely to be sampled and to reappear afterwards, contributing proportionally to $m_j^2$.

\subsubsection*{Implementation detail}
The $k$ sampled positions are stored in sorted order, so the update loop can break
early as soon as it encounters a position beyond the current stream index.

On average only half the variables are active at any point in the stream, so the
expected number of iterations per update call is $k/2$. I used \texttt{AMS\_NUM\_VARS = 100}.

\subsubsection*{Results}

\begin{table}[H]
\centering
\begin{tabular}{rrr}
\toprule
$k$ (variables) & Estimate & Relative error \\
\midrule
 10 & 9,246,760,000 & 4.3\% \\
 25 & 9,891,080,000 & 2.3\% \\
 50 & 9,677,872,000 & 0.1\% \\
100 & 9,214,792,000 & 4.7\% \\
200 & 9,716,568,000 & 0.5\% \\
\bottomrule
\end{tabular}
\caption{AMS on the subsample (200,000 comments, \texttt{SEED=1}, exact $F_2^* = 9{,}667{,}184{,}406$).}
\label{tab:ams_subsample}
\end{table}

\begin{table}[H]
\centering
\begin{tabular}{rrr}
\toprule
$k$ (variables) & Estimate & Relative error \\
\midrule
 10 & 4,094,686,303,929 & 30.8\% \\
 25 & 5,649,804,521,452 &  4.5\% \\
 50 & 6,104,101,436,235 &  3.2\% \\
100 & 5,234,470,416,714 & 11.5\% \\
200 & 5,269,086,030,059 & 10.9\% \\
\bottomrule
\end{tabular}
\caption{AMS on the full dataset (4,986,461 comments, \texttt{SEED=1}, exact $F_2^* = 5{,}916{,}853{,}781{,}669$).}
\label{tab:ams}
\end{table}

The estimator is non-monotone in $k$ for a single run, consistent with its
theoretical variance of $O(F_2^2/k)$.
On the subsample errors stay below 5\% for all $k$; on the full dataset errors
reach 10--30\% even at large $k$, reflecting the higher $F_2$ and wider variance
of the larger stream.

\subsection{Bloom Filter (Set Membership)}

Given the set $S$ of Opinion article IDs, the goal is to answer membership queries
``is article $x \in S$?'' in $O(1)$ time, tolerating a small rate of false positives.

\subsubsection*{Data structure}
A Bloom filter consists of a bit array of $m$ bits (initially all zero) and $k$
independent hash functions $h_1, \ldots, h_k : U \to [0, m)$.
\begin{itemize}
  \item \textbf{Insertion of $x$}: set bits $h_1(x), \ldots, h_k(x)$ to 1.
  \item \textbf{Query for $x$}: return \textsc{true} iff all $k$ bits are set.
\end{itemize}
False negatives are impossible; false positives occur when all $k$ hash positions of
a non-member element happen to be set by other insertions.
The theoretical false positive rate (FPR) for $n$ inserted elements is:
\[
  \mathrm{FPR}(m,k,n) = \left(1 - e^{-kn/m}\right)^k
\]
The optimal number of hash functions for a given $m/n$ ratio is
$k^* = (m/n)\ln 2$.

\subsubsection*{Bit-array size}
The number of hash functions is fixed at \texttt{BF\_K\_HASHES} $= k = 7$.
This value is an upper bound on $k^*$ for all targets tested with the full dataset
($n = 2{,}272$): the most stringent target of $p = 1\%$ gives
$k^* \approx 6.65$, and larger targets give smaller $k^*$.

Rather than fixing $m$ as a global constant, the filter computes it automatically
from $n$, $k$, and a target FPR $p$ by solving $\mathrm{FPR}(m,k,n) = p$
for $m$ with $k$ fixed:
\[
  m = \left\lceil \frac{-k\, n}{\ln\!\left(1 - p^{1/k}\right)} \right\rceil
\]
The default target is \texttt{BF\_TARGET\_FPR = 0.01} (1\%), giving
$m \approx 21{,}796$ bits for the 2,272 full-dataset Opinion articles.

\subsubsection*{Hash functions}
Each of the $k$ hash functions is implemented as:
\[
  h_i(x) = \mathrm{MD5}(x \,\|\, \mathrm{salt}_i) \bmod m
\]
where $\mathrm{salt}_i$ is a distinct random 32-bit integer string generated at
construction time; MD5 works well here since article IDs are URI strings.

\subsubsection*{Results}

\begin{table}[H]
\centering
\begin{tabular}{rrrr}
\toprule
Target FPR & $m$ (bits) & FPR empirical & FPR theoretical \\
\midrule
50\% &  6,735 & 50.40\% & 50.00\% \\
20\% & 10,049 & 19.44\% & 19.99\% \\
10\% & 12,505 &  9.67\% & 10.00\% \\
 5\% & 15,074 &  4.99\% &  5.00\% \\
 1\% & 21,796 &  1.01\% &  1.00\% \\
\bottomrule
\end{tabular}
\caption{Bloom filter on the full dataset ($n=2{,}272$ Opinion articles, $k=7$,
  14,515 negative queries, \texttt{SEED=1}).}
\label{tab:bloom}
\end{table}

\begin{table}[H]
\centering
\begin{tabular}{rrrr}
\toprule
Target FPR & $m$ (bits) & FPR empirical & FPR theoretical \\
\midrule
50\% &   338 & 60.18\% & 49.98\% \\
20\% &   505 & 19.30\% & 19.91\% \\
10\% &   628 &  9.65\% &  9.97\% \\
 5\% &   757 &  4.91\% &  4.98\% \\
 1\% & 1,094 &  0.70\% &  1.00\% \\
\bottomrule
\end{tabular}
\caption{Bloom filter on the subsample ($n=114$ Opinion articles, $k=7$,
  570 negative queries, \texttt{SEED=1}).}
\label{tab:bloom_subsample}
\end{table}

The theoretical FPR matches the target across all settings.
On the full dataset the empirical FPR tracks the target closely for all values of
$p$; on the subsample the 50\% target shows a larger deviation (60\% vs.\ 50\%)
due to the high variance of the observed proportion when both $n$ (114 inserted
elements) and the number of negative queries (570) are small.

\section{Scalability}

\begin{table}[H]
\centering
\begin{tabular}{llll}
\toprule
Algorithm & Memory (algorithm state) & Time per element  \\
\midrule
Flajolet-Martin  & $O(H)$ & $O(H)$ \\
AMS              & $O(k)$ & $O(k)$ \\
Bloom filter     & $O(m)$ & $O(k_{\mathrm{BF}})$ \\
\bottomrule
\end{tabular}
\caption{Complexity of the three algorithms. $H$ = number of hash functions (FM),
$k$ = number of AMS variables, $m$ = number of Bloom filter bits,
$k_{\mathrm{BF}}$ = number of Bloom filter hash functions. }
\label{tab:scalability}
\end{table}

The memory used by the algorithm \emph{state} is constant with respect to the stream
length $n$: FM stores $H = 64$ integers, AMS stores $k = 100$ counters and element
labels, and the Bloom filter stores an $m$-bit array where $m$ is determined by $n$
and the target FPR (e.g.\ $m \approx 21{,}796$ bits $\approx$2.7~KB for $n=2{,}272$
and target 1\%).


\begin{thebibliography}{9}

\bibitem{mmds}
J.~Leskovec, A.~Rajaraman, and J.~D.~Ullman,
\textit{Mining of Massive Datasets}, 3rd ed.,
Cambridge University Press, 2020.
\url{http://www.mmds.org/}

\bibitem{kaggle}
B.~Awad,
\textit{New York Times Articles \& Comments (2020)},
Kaggle, 2020.
\url{https://www.kaggle.com/datasets/benjaminawd/new-york-times-articles-comments-2020}

\end{thebibliography}

\section*{Declaration}

I declare that this material, which I now submit for assessment, is entirely my own
work and has not been taken from the work of others, save and to the extent that such
work has been cited and acknowledged within the text of my work. I understand that
plagiarism, collusion, and copying are grave and serious offences in the university
and accept the penalties that would be imposed should I engage in plagiarism,
collusion or copying. This assignment, or any part of it, has not been previously
submitted by me or any other person for assessment on this or any other course of
study. No generative AI tool has been used to write the code or the report content.

\end{document}
